% Copyright (c) 2026 Eugene Joseph M. Ragasa.
% SPDX-License-Identifier: Apache-2.0
% Author: Eugene Joseph M. Ragasa
% Formal proof completed: 2026-10-27
\documentclass[11pt]{article}
\usepackage[T1]{fontenc}\usepackage{amsmath,amssymb}\usepackage[margin=1in]{geometry}\usepackage[hidelinks]{hyperref}
\title{Proof Dossier: Character-Phase Fourier Mode Shift}\author{Eugene Joseph M. Ragasa}\date{Formal proof completed: 2026-10-27}
\begin{document}\maketitle
\section{Identity and status}
ID \texttt{BG-FOURIER-THM-048}; status \texttt{proved}. Lean declaration:
\par\noindent{\tiny\path{ProjectKoios.BandGauge.Fourier.}}\\
{\tiny\path{finiteFourierCoefficient_characterPhaseMultiply}}.
\section{Statement and proof}
For the normalized conjugated-character forward transform, multiplying a vector-valued
field by character $p$ sends the coefficient at $m$ to the original coefficient at
$m-p$. Pointwise character subtraction and unit-modulus conjugation reduce each summand.
\section{Boundary}
This is a state/vector action and one pure-character mode shift, not scalar operator
conjugation or a general pointwise-gauge convolution theorem. Kernel axioms are
\texttt{propext}, \texttt{Classical.choice}, and \texttt{Quot.sound}; no admissions.
Checked with Lean \texttt{4.35.0} and Mathlib revision
\texttt{56a005855ea17e0b0958f73ec801ae140b0a1679}. AI-assisted and kernel-checked;
no independent expert review recorded.
\end{document}
